\section{Worum es geht}

\textbf{[SETZUNG]} Dieses Dokument entwickelt die vollständige Skalenhierarchie der
FFGFT: von der einzigen geometrischen Konstante $\xi$ zu den charakteristischen
Energieskalen, zu den Leptonmassen, zur Feinstrukturkonstante $\alpha$ und zu $E_0$
als logarithmischer Mitte. Der Kern: alle diese Größen folgen aus $\xi$ -- aber die
Verbindung ist nicht immer durch direkte Kürzung herzustellen. Verhältnisse
(Massenverhältnisse) sind korrekturfrei; Absolutwerte (auch $E_0$ und $\alpha$) brauchen $K_\text{frak}$.
\textit{(Korrigiert am 30.~Sept.~2026: vorher waren Massenverhältnisse und $E_0$ als korrekturfrei genannt -- $E_0 = \sqrt{m_e m_\mu}$ ist ein Absolutwert; der korrigierte Wert ist $E_0/\sqrt{K_\text{frak}} = 7{,}398$~MeV (Dok.~190, R72).)}

\section{Der Ausgangspunkt: eine einzige Konstante}

\textbf{[B]} Die FFGFT beginnt mit einer einzigen dimensionslosen Konstante:
\[
  \xi = \frac{4}{3}\times10^{-4} = 1{,}333\ldots\times10^{-4}.
\]
Aus ihr folgen in natürlichen Einheiten ($\hbar=c=\ell_P=1$) alle fundamentalen Skalen.

\textbf{[B]} \textbf{Charakteristische T0-Skalen:}
\[
  L_0 = \xi\cdot\ell_P \approx 1{,}333\times10^{-4}\,\ell_P,
  \qquad
  E_0^\text{Skala} = L_0^{-1} = \xi^{-1}\cdot E_P.
\]

\section{Leptonmassen aus Geometrie}

\textbf{[K]} Die Massenformeln (A100) in kompakter Form:
\[
  m_e = c_e\cdot\xi^{5/2},
  \quad
  m_\mu = c_\mu\cdot\xi^2,
  \quad
  m_\tau = c_\tau\cdot\xi^{3/2},
\]
wobei die Koeffizienten geometrischen Ursprung haben:
\[
  c_e = \frac{3\sqrt3}{2\pi\alpha^{1/2}},
  \qquad
  c_\mu = \frac{9}{4\pi\alpha},
  \qquad
  c_\tau = \frac{27\sqrt3}{8\pi\alpha^{3/2}}.
\]

\textbf{[K]} \textbf{Wichtig: scheinbar einfache Brüche.} In der Kurzform
erscheinen die Koeffizienten als rationale Zahlen ($r_e=4/3$, $r_\mu=16/5$,
$r_\tau=25/9$ aus A100). Diese Brüche sind \emph{nicht} willkürlich und dürfen
nicht direkt gekürzt werden: sie kodieren Raumgeometrie ($\pi$, $\sqrt3$) und
elektromagnetische Kopplung ($\alpha$). Beide Darstellungen (Yukawa-Kurzform
und erweiterte Form) sind dieselbe Theorie in zwei Schreibweisen (A130).
\textit{(Vermerk vom 30.~Sept.~2026: Die erweiterte Form ist mit A100 nicht äquivalent. Mit $c_e = 9{,}681$, $c_\mu = 98{,}14$, $c_\tau = 2985$ folgt $m_\mu/m_e = (c_\mu/c_e)\,\xi^{-1/2} = 878$ (gemessen $206{,}77$) und $m_\tau/m_\mu = (c_\tau/c_\mu)\,\xi^{-1/2} = 2634$ (gemessen $16{,}82$); die Exponentendifferenz $\tau/\mu$ ist hier $-1/2$ statt $-1/3$ wie in A100 ($p = 3/2, 1, 2/3$). Die Massenverhältnisse werden von dieser Form nicht reproduziert; maßgeblich ist die Leiter $m_i = r_i\,\xi^{p_i}\,v$ aus A100.)}

\section{E₀ als logarithmische Mitte}

\textbf{[B]} Die charakteristische Energie $E_0$ ist das geometrische Mittel der
Elektron- und Myonmasse:
\[
  \boxed{E_0 = \sqrt{m_e\cdot m_\mu}}.
\]
Logarithmisch: $\log E_0 = \tfrac12[\log m_e+\log m_\mu]$ -- $E_0$ liegt exakt in
der Mitte zwischen $m_e$ und $m_\mu$ auf logarithmischer Skala. Numerisch:
$E_0 = \sqrt{0{,}511\times105{,}658}\,\text{MeV}\approx7{,}35\,\text{MeV}$ (geometrisches Mittel). Der fraktal korrigierte Wert (A130, Weg~2) ist $E_0 = 7{,}348/\sqrt{K_\text{frak}}\approx7{,}398\,\text{MeV}$; diese Differenz ist die Stelle, wo $K_\text{frak}$ eingeht. \textit{(Korrigiert am 30.~Sept.~2026: vorher \enquote{Der Wert aus dem $\xi$-Weg (A130)} -- dieser Weg, $E_0^2 = 4\sqrt2\,m_\mu/\xi^{p}$, ist in A130 entfernt.)}

\textbf{[B]} Aus den Massenformeln:
\[
  E_0 = \sqrt{c_ec_\mu}\cdot\xi^{9/4}.
\]
$E_0$ ist \emph{nicht} korrekturfrei: im geometrischen Mittel kürzt sich
$K_\text{frak}$ nicht weg. Der fraktal korrigierte Wert ist
$E_0 = E_0^\text{bare}/\sqrt{K_\text{frak}} = 7{,}348/\sqrt{0{,}98667} = 7{,}397\approx7{,}398\,\text{MeV}$
(R72); die Wurzel folgt daraus, dass $E_0$ quadratisch in $\alpha$ eingeht.
\textit{(Korrigiert am 30.~Sept.~2026: vorher \emph{korrekturfrei} und $E_0^\text{exp}=K_\text{frak}\cdot E_0^\text{bare}$ -- $0{,}98667\times7{,}348 = 7{,}250$~MeV entfernt sich von $7{,}398$~MeV; die Korpusbeziehung ist $E_0^\text{bare}/\sqrt{K_\text{frak}}$ (Dok.~190, R72).)}

\section{Die Feinstrukturkonstante aus $\xi$}

\textbf{[K]} Die fundamentale Beziehung:
\[
  \alpha = \xi\cdot E_0^2 = c_ec_\mu\cdot\xi^{11/2}.
\]
Einsetzen der geometrischen Koeffizienten:
\[
  \alpha^{5/2} = \frac{27\sqrt3}{8\pi^2}\cdot\xi^{11/2},
  \qquad\text{also}\qquad
  \alpha = \Bigl(\frac{27\sqrt3}{8\pi^2}\Bigr)^{2/5}\cdot\xi^{11/5}.
\]
Mit fraktaler Korrektur (A040):
\[
  \alpha = \Bigl(\frac{27\sqrt3}{8\pi^2}\Bigr)^{2/5}\cdot\xi^{11/5}\cdot K_\text{frak},
  \qquad K_\text{frak}=1-100\xi\approx0{,}9867.
\]
Numerisch ergibt diese geschlossene Form mit $\xi = \tfrac43\times10^{-4}$ allerdings
$\alpha = 0{,}811\times2{,}98\times10^{-9}\times0{,}9867 = 2{,}4\times10^{-9}$
($\alpha^{-1}\approx4{,}2\times10^{8}$), nicht $1/137$. Der Wert $\alpha^{-1} = 137{,}035$
($0{,}0005\,\%$) folgt aus $\alpha = \xi E_0^2$ mit $E_0 = 7{,}398$~MeV (A130, R72).
\textit{(Korrigiert am 30.~Sept.~2026: vorher \enquote{Numerisch: $\alpha^{-1}\approx137{,}036$ -- Übereinstimmung auf $0{,}0005\,\%$} für die geschlossene Form -- $(27\sqrt3/8\pi^2)^{2/5} = 0{,}811$, $\xi^{11/5} = 2{,}98\times10^{-9}$; die Kette $\alpha = c_ec_\mu\xi^{11/2}$ reproduziert $\alpha$ nicht, vgl.\ Dok.~070.)}

\section{Warum $\alpha$ nicht direkt aus $\xi^{11/2}$ folgt ohne Einheitenkonvention}

\textbf{[S]} $\alpha\propto\xi^{11/2}$ bedeutet: wenn sich $\xi$ änderte, änderte
sich $\alpha$ mit Exponent $11/2=5{,}5$. Das ist eine starke Kopplung. Die Koeffizienten
$c_e,c_\mu$ sind in natürlichen Einheiten von Größenordnung $10^{19}$ -- das ist kein
freier Parameter, sondern der Umrechnungsfaktor zwischen Planck-Skala und MeV. In
natürlichen Einheiten (Planck-Energien) werden $c_e\sim9{,}67$ und $c_\mu\sim98{,}1$ --
Größenordnung 1 bis 100, geometrisch verständlich.
\textit{(Vermerk vom 30.~Sept.~2026: $c_e\approx9{,}68$ und $c_\mu\approx98{,}1$ sind die aus $\alpha$ berechneten Zahlen der Koeffizientenformeln, keine Werte in Planck-Einheiten; in Planck-Einheiten wäre $m_e/\xi^{5/2} = 4{,}19\times10^{-23}/2{,}05\times10^{-10}\approx2{,}0\times10^{-13}$.)}

\section{Verhältnisse sind korrekturfrei}

\textbf{[B]} Das Massenverhältnis:
\[
  \frac{m_\mu}{m_e}
  = \frac{K_\text{frak}\cdot m_\mu^\text{bare}}{K_\text{frak}\cdot m_e^\text{bare}}
  = \frac{m_\mu^\text{bare}}{m_e^\text{bare}}.
\]
$K_\text{frak}$ kürzt sich heraus. Dasselbe gilt für die Koide-Verhältnisse (A110)
und die $g-2$-Verhältnisse (A138). Das ist kein Zufall: multiplikative Korrekturen
kürzen sich in Verhältnissen immer heraus.

\textbf{[K]} Absolutwerte dagegen brauchen $K_\text{frak}$: ohne Korrektur liegt
$\alpha^{-1}\approx138{,}9$ statt $137{,}04$ -- $1{,}4\,\%$ Abweichung. Mit
$K_\text{frak}$: $0{,}98667\times138{,}91 = 137{,}06$ ($+0{,}017\,\%$).
\textit{(Korrigiert am 30.~Sept.~2026: vorher \enquote{Mit $K_\text{frak}$: exakt} -- $K_\text{frak}/(\xi\cdot7{,}3479^2) = 137{,}059$ gegen $137{,}036$.)} Damit ist $K_\text{frak}$ keine Nachkorrektur, sondern
Bestandteil der Struktur (A040).

\section{Die vollständige Ableitungskette}

\textbf{[B]} Alle fundamentalen Größen aus $\xi$ in einer Kette:
\begin{align}
  \xi
  &\;\xrightarrow{\text{Geometrie}}\; r_0,E_0^\text{Skala}
   \;\xrightarrow{\text{Quantenzahlen}}\; m_e,m_\mu,m_\tau \notag\\
  &\;\xrightarrow{\text{geom. Mittel}}\; E_0
   \;\xrightarrow{\xi\cdot E_0^2}\; \alpha
   \;\xrightarrow{\text{K-Korr.}}\; \alpha_\text{exp}. \notag
\end{align}
Zusätzlich: $\xi\to L_0\to G$ (Gravitationskonstante, A142) und $\xi\to L_0$ (A015).

\section{Prüfskript}

\begin{itemize}[leftmargin=2em,itemsep=2pt,topsep=2pt]
  \item Numerisch: $E_0=\sqrt{m_em_\mu}$, $\alpha=\xi E_0^2$, Verhältnisse
    korrekturfrei vs.\ Absolutwerte mit $K_\text{frak}$:
    \texttt{python/A\_Serie\_Skripte/a261\_skalenhierarchie.py}
\end{itemize}

\section{Was offen bleibt}

\textbf{Die Koeffizienten $c_e,c_\mu,c_\tau$ sind nicht vollständig unabhängig
hergeleitet.} Sie enthalten $\alpha$ explizit -- und $\alpha$ wiederum folgt aus den
Massen. Die Zirkularität ist in A130 als methodisches Problem benannt; die T0-Lösung
(gemeinsamer geometrischer Ursprung) ist eine Behauptung, kein Beweis.

\textbf{Die Quark-Koeffizienten fehlen.} Die $c_i$-Struktur für Quarks ist nicht
auf demselben Niveau ausgearbeitet wie für Leptonen (A150).

\section{Ersetzt}

Dieses Dokument nimmt auf: Dok.~070 (mathematische Struktur der T0-Theorie:
zirkuläre Verhältnisse, Skalenhierarchie aus $\xi$, Leptonmassen, $\alpha$-Herleitung,
$E_0$ als logarithmische Mitte, geometrische Konstante $C$, warum Verhältnisse
korrekturfrei sind). Dok.~093 ($\beta$-Parameter-Ableitung, Selbstkonsistenzbedingung).

\section{Verwendung von KI-Werkzeugen}

Dieses Dokument ist unter Verwendung KI-gestützter Werkzeuge entstanden. Verantwortung
für Inhalt und Fehler liegt beim Autor. Wortlaut der Regeln in A010.
